\documentclass[12pt,a4paper]{article}

\usepackage{fontspec}
\setmainfont{Liberation Serif}
\setsansfont{Liberation Sans}
\setmonofont{Liberation Mono}
\linespread{1.0}
\fontsize{14pt}{17pt}\selectfont

\usepackage{geometry}
\geometry{left=3cm,right=1.5cm,top=2cm,bottom=2cm}

\usepackage{indentfirst}
\setlength{\parindent}{1.25cm}

\usepackage{array}
\usepackage{tabularx}

\usepackage{fancyhdr}
\pagestyle{fancy}
\fancyhf{}
\fancyfoot[C]{\thepage}
\renewcommand{\headrulewidth}{0pt}

\begin{document}

% Титульная страница
\thispagestyle{empty}
\begin{center}
ГОСУДАРСТВЕННОЕ УЧРЕЖДЕНИЕ ОБРАЗОВАНИЯ\\
<<МИНСКИЙ ГОСУДАРСТВЕННЫЙ КОЛЛЕДЖ\\
ЭЛЕКТРОНИКИ>>
\end{center}

\vspace{3cm}

\begin{center}
{\bfseries ОТЧЁТ\\
ПО ФУНКЦИОНАЛЬНОМУ ТЕСТИРОВАНИЮ\\
ПРОЕКТА <<GoldPC>>}
\end{center}

\vspace{2cm}

\begin{center}
2026
\end{center}

\vfill

\newpage

\section*{СОДЕРЖАНИЕ}

1.~Общие сведения о проекте\\
2.~Объект тестирования\\
3.~Методы и средства тестирования\\
4.~Результаты функционального тестирования\\
5.~Выводы

\newpage

\section{1. Общие сведения о проекте}

Проект <<GoldPC>> представляет собой интернет-магазин персональных компьютеров с функционалом конфигуратора сборки ПК, каталогом товаров, системой заказов, сервисным обслуживанием и гарантийным обслуживанием. Система построена на архитектуре микросервисов и включает следующие компоненты: AuthService, CatalogService, OrdersService, ServicesService, WarrantyService, PCBuilderService, GoldPC.Api.

Стек технологий: ASP.NET Core 8 (бэкенд), React 19 + Vite (фронтенд), PostgreSQL (база данных), gRPC (межсервисное взаимодействие), Redis (кэширование).

\section{2. Объект тестирования}

Объектом тестирования является бэкендная часть системы, реализующая бизнес-логику приложения. Функциональное тестирование проводилось на уровне модульных тестов для каждого микросервиса.

\section{3. Методы и средства тестирования}

Для функционального тестирования использовались: модульные тесты (xUnit, FluentAssertions, Moq), интеграционные тесты (WebApplicationFactory), контрактные тесты (PactNet), in-memory база данных (EF Core InMemory).

\section{4. Результаты функционального тестирования}

В таблице представлены результаты 7 основных функциональных тестов.

\begin{center}
\begin{tabularx}{\textwidth}{|c|X|c|}
\hline
№ & Тест & Результат \\
\hline
1 & Register\_ValidRequest (AuthService) & Пройден \\
\hline
2 & GetProductById (CatalogService) & Пройден \\
\hline
3 & CreateOrder (OrdersService) & Пройден \\
\hline
4 & CreateServiceRequest (ServicesService) & Пройден \\
\hline
5 & CreateWarrantyClaim (WarrantyService) & Пройден \\
\hline
6 & CheckCompatibility (PCBuilderService) & Пройден \\
\hline
7 & VerifyPact (ContractTests) & Пройден \\
\hline
\end{tabularx}
\end{center}

Всего тестов: 7. Пройдено: 7. Не пройдено: 0.

\section{5. Выводы}

В ходе функционального тестирования проекта <<GoldPC>> выполнена проверка основных сценариев работы системы. Все 7 тестов прошли успешно, что подтверждает корректность реализации бизнес-логики микросервисов.

\end{document}
